\section{Binary lift forcing}

Model the visible orientation cycle by vertices \(0,1,2,3\) and a
successor \(i\mapsto i+1\pmod 4\). A deterministic two-sheet lift is
specified by four edge voltages
\[
\epsilon_0,\epsilon_1,\epsilon_2,\epsilon_3\in C_2.
\]
Its circuit holonomy is
\[
\omega=\epsilon_0+\epsilon_1+\epsilon_2+\epsilon_3\in C_2.
\]

Changing the sheet labels over individual visible vertices is a gauge
transformation. Gauge transformation changes the individual edge
voltages but preserves \(\omega\).

\begin{theorem}[Binary lift classification]
Every deterministic two-sheet lift of the visible four-cycle belongs
to one of two gauge classes.
\begin{enumerate}
\item If \(\omega=0\), the lift is two disjoint four-cycles.
\item If \(\omega=1\), the lift is one connected eight-cycle.
\end{enumerate}
\end{theorem}

\begin{proof}
After four lifted successor steps, the visible vertex returns and the
sheet changes by \(\omega\). Trivial holonomy therefore closes each
sheet separately after four steps. Nontrivial holonomy exchanges the
sheets after four steps and returns to the initial lifted state only
after eight. Gauge normalization moves all nonzero voltage to one
chosen edge, proving uniqueness within each holonomy class.
\end{proof}

Hence a nontrivial binary full-circuit receipt uniquely forces the
connected \(C_8\) lift, up to fiber-preserving gauge and isomorphism.
Its deck involution is the four-step element \(H^4=\tau\).
